\chapter*{Preface}

\textit{Wild Luck} is a companion to \textit{Luck}.  The first book defined
luck and computed it.  This one uses it.

The bet is that one definition --- $L(x) = |\Omega(x)| + \tfrac{1}{2}|\omega(x)|$,
the probability of an outcome at least as common as $x$, with ties
half-counted --- is enough machinery to do most of the everyday statistical
work that real domains demand: rejecting a forecast as miscalibrated,
deciding when a filter has lost track of its system, choosing between
competing models, recognizing when ``rare'' has crossed into ``too rare to
believe.''  Each chapter pulls one of two levers on the same statistic.

\paragraph{The Read lever} ``Given a model and an observation, what is
$\ell(x)$, and what should I do about it?''  The answer drives one-shot
interpretation, anomaly detection, hypothesis testing, and any decision
where someone has to act on a single event.

\paragraph{The Drive lever} ``Given many observations, find a model where
$\ell$ is uniform.''  The answer drives parameter fitting, state
estimation, drift tracking, and model selection.  Same statistic, opposite
direction.

Every chapter that follows is a worked example of one of these two moves
in a specific domain.  The unification is the point: someone who has
internalized the Read/Drive duality can carry it from weather forecasting
to Kalman filtering to clinical decision-making without learning a new
hammer for every nail.

\paragraph{What this book is not.}  It is not a survey of optimal methods.
The Kalman filter has decades of optimization literature; classical
hypothesis testing has Neyman--Pearson; model selection has AIC/BIC.  Each
chapter compares luck-based methods to those alternatives, honestly.  Luck
wins on \textit{transferability} --- one statistic, one diagnostic,
graded against one noise floor --- not on raw efficiency.  When you know
the alternative, the Neyman--Pearson likelihood-ratio test will reject
faster.  When you do not, $\ell$ is what is left, and that is most of the
real cases.

\paragraph{What this book assumes.}  The reader has seen \textit{Luck} or
is willing to consult it.  Each chapter opens with an
\textit{Ingredients} box pointing to the specific theorems, definitions,
and recipes drawn from the first book.  The chapters do not re-derive.
That keeps this book a catalog, not a textbook --- read any chapter
independently, in any order, after Chapter~\ref{ch:coin}.

\paragraph{Notation.}  $L$ is the analytic luck $|\Omega| + \tfrac{1}{2}
|\omega|$.  $\ell$ is the count-based estimator from a sample.  $C =
\log_2(\ell/(1-\ell))$ is the \emph{coin-luck}: how many coin-flips ahead
of (positive) or behind (negative) typical, an outcome lies.  These are
the running characters of every chapter.
